\documentclass[11pt,letterpaper]{article}
\usepackage[numbers,sort&compress]{natbib}
\usepackage{angel-publication}
\PublicationRunningTitle{Manuscript title}
\PublicationRunningAuthor{Author name}
\PublicationStatus{House-style template v1.0 - 24 October 2026}
\hypersetup{pdftitle={Manuscript title - house-style specimen},pdfauthor={Author name}}
\title{Manuscript Title\\[4pt]{\large\normalfont\color{black}Optional descriptive subtitle}}
\author{Author Name\\\small Affiliation or independent researcher\\
  \small\href{mailto:author@example.org}{author@example.org}}
\date{24 October 2026\\[6pt]
  \PublicationLinks{https://example.org/report/}{https://github.com/example/project}}

\begin{document}
\maketitle
\begin{abstract}
This specimen demonstrates the shared presentation standard. Replace the example
material with the manuscript's actual scientific content. State the question,
method, result and limits directly; distinguish a proved statement from an
experimental observation. The examples below are elementary formatting tests,
not research results.
\end{abstract}

\section{Introduction}
\label{sec:introduction}
Use a single column, 11-point body type, Latin Modern text and mathematics, and
one-inch margins on US Letter paper. All document-specific notation belongs in
this source file. The style package owns typography and page layout. A citation
uses the manuscript's existing reference data, for example \citet{lamport}.

\section{Definitions and results}
\label{sec:results}
\begin{definition}[Finite arithmetic mean]
For real numbers $x_1,\ldots,x_n$ with $n\geq1$, define
\begin{equation}
  \overline{x}=\frac{1}{n}\sum_{i=1}^{n}x_i.
  \label{eq:mean}
\end{equation}
\end{definition}

\begin{lemma}
\label{lem:constant}
If $x_i=c$ for every $i\in\{1,\ldots,n\}$, then $\overline{x}=c$.
\end{lemma}
\begin{proof}
Substitute $x_i=c$ into \eqref{eq:mean}. The sum is $nc$, and $n\ne0$.
\end{proof}

\begin{theorem}[Translation of a mean]
\label{thm:translation}
If $y_i=x_i+a$ for every $i$, then $\overline{y}=\overline{x}+a$.
\end{theorem}
\begin{proof}
The algebra can be displayed without overlong single lines:
\begin{align*}
\overline{y}
  &=\frac{1}{n}\sum_{i=1}^{n}(x_i+a)\\
  &=\frac{1}{n}\sum_{i=1}^{n}x_i+\frac{na}{n}
   =\overline{x}+a.\qedhere
\end{align*}
\end{proof}

\begin{corollary}
Centering the $x_i$ by subtracting $\overline{x}$ gives a mean of zero.
\end{corollary}
\begin{proof}
Apply \cref{thm:translation} with $a=-\overline{x}$.
\end{proof}
\begin{remark}
Theorem-like statements share a counter within each section. Proofs use the AMS
proof environment, including a correctly placed end-of-proof mark.
\end{remark}

\section{Figures, tables and references}
\label{sec:presentation}
\begin{figure}[htbp]
  \centering
  \setlength{\unitlength}{1cm}
  \begin{picture}(12,1.5)
    \put(0,.2){\framebox(3,1){Inputs}}
    \put(3,.7){\vector(1,0){1.5}}
    \put(4.5,.2){\framebox(3,1){Declared method}}
    \put(7.5,.7){\vector(1,0){1.5}}
    \put(9,.2){\framebox(3,1){Measured result}}
  \end{picture}
  \caption{A vector diagram with labels at a readable size. The caption explains
  the relationship and the role of each object.}
  \label{fig:method}
\end{figure}

\begin{table}[htbp]
  \centering
  \caption{Formatting examples only. Numeric columns are right aligned, use
  consistent precision, and declare the quantity or unit in the header.}
  \label{tab:example}
  \begin{tabular}{@{}lrrr@{}}
    \toprule
    Case & Count & Mean (units) & Difference (units)\\
    \midrule
    A & 8 & 12.30 & $+0.20$\\
    B & 16 & 9.75 & $-0.35$\\
    \bottomrule
  \end{tabular}
\end{table}

Refer to objects by labels rather than hard-coded numbers:
\cref{sec:results,lem:constant,fig:method,tab:example}; use
\verb|\eqref| for a parenthesized equation reference. Keep units and precision
consistent and preserve uncertainty, sample sizes and denominators. Use
\verb|booktabs| rules, with no vertical rules. Put table captions above tables
and figure captions below figures.

Use \verb|[htbp]| for normal floats and explicit \verb|\FloatBarrier| only when
the scientific reading order requires it. Oversized diagrams can use
\verb|publicationdiagram|; this wrapper shrinks to the available width without
changing labels or data. Reflow long equations rather than shrinking them.
URLs remain live and wrap, for example
\url{https://www.latex-project.org/help/documentation/}.

\subsection{Resource links follow availability}
The front matter uses Overview, Code and Lean in that order, omitting unavailable
items. Empty arguments produce no link label or separator. These specimen rows
exercise the optional formalization link and the absent-code cases:
\begin{center}
\begin{tabular}{@{}ll@{}}
Overview, code and formalization &
\PublicationLinks[https://example.org/formalization/]{https://example.org/report/}{https://github.com/example/project}\\[4pt]
Overview and formalization &
\PublicationLinks[https://example.org/formalization/]{https://example.org/report/}{}\\[4pt]
Overview only &
\PublicationLinks{https://example.org/report/}{}
\end{tabular}
\end{center}

\FloatBarrier
\section{Bibliography and build acceptance}
For a new manuscript use numbered natbib citations and \verb|plainnat|, as this
specimen does. Existing manuscripts retain their complete bibliography and
citation semantics when adopted. Compile until all references and citations
resolve. Inspect the final PDF's pages, not only the compiler log.

\bibliographystyle{plainnat}
\bibliography{references}
\end{document}
